\begin{lemma}
For every \(n\) for which the real reduced count and real independent-set scale
are nonnegative,
\[
M(n)\le \noderef{TwoBiteReducedVertexCount}(n)<M(n)+1
\]
and
\[
\noderef{TwoBiteIndependenceScale}(\kappa,n)
\le K_\kappa(n)
<\noderef{TwoBiteIndependenceScale}(\kappa,n)+1,
\]
where \(M=\noderef{TwoBiteNaturalReducedVertexCount}\) and
\(K_\kappa=\noderef{TwoBiteNaturalIndependenceScale}\).
\end{lemma}
\begin{proof}
Fix \(\kappa\) and \(n\), and assume
\[
0\le \noderef{TwoBiteReducedVertexCount}(n),
\qquad
0\le \noderef{TwoBiteIndependenceScale}(\kappa,n).
\]
Put
\[
m_R=\noderef{TwoBiteReducedVertexCount}(n),\qquad
k_R=\noderef{TwoBiteIndependenceScale}(\kappa,n),
\]
\[
M(n)=\noderef{TwoBiteNaturalReducedVertexCount}(n),
\qquad
K_\kappa(n)=\noderef{TwoBiteNaturalIndependenceScale}(\kappa,n).
\]

Unfolding \(\noderef{TwoBiteNaturalReducedVertexCount}\), the integer
\(M(n)\) is \(\lfloor m_R\rfloor\), the natural floor of \(m_R\).  The first
nonnegativity hypothesis says \(m_R\ge0\), so the natural floor is governed by
the usual real floor inequalities:
\[
\lfloor m_R\rfloor\le m_R
\qquad\text{and}\qquad
m_R<\lfloor m_R\rfloor+1.
\]
Substituting \(M(n)=\lfloor m_R\rfloor\) gives the first two conjuncts, in the
same order as the Lean statement:
\[
(M(n):\mathbb R)\le m_R
\]
and
\[
m_R<(M(n):\mathbb R)+1.
\]

Unfolding \(\noderef{TwoBiteNaturalIndependenceScale}\), the integer
\(K_\kappa(n)\) is \(\lceil k_R\rceil\), the natural ceiling of \(k_R\).  The
second nonnegativity hypothesis says \(k_R\ge0\), so the natural ceiling
inequalities apply in real-valued form:
\[
k_R\le \lceil k_R\rceil
\qquad\text{and}\qquad
\lceil k_R\rceil<k_R+1.
\]
Substituting \(K_\kappa(n)=\lceil k_R\rceil\) gives the remaining two conjuncts,
again in the order of the Lean statement:
\[
k_R\le (K_\kappa(n):\mathbb R)
\]
and
\[
(K_\kappa(n):\mathbb R)<k_R+1.
\]
Replacing \(m_R\) and \(k_R\) by their definitions yields exactly the four
claimed inequalities.
\end{proof}
